% Part: first-order-logic
% Chapter: natural-deduction
% Section: rules-and-proofs

\documentclass[../../../include/open-logic-section]{subfiles}

\begin{document}

\iftag{FOL}
      {\olfileid{fol}{ntd}{rul}}
      {\olfileid{pl}{ntd}{rul}}

\olsection{નિયમો અને \usetoken{P}{derivation}}

\begin{explain}
સ્વાભાવિક નિગમનનાં તંત્રો ગણિતીય સાબિતીમાં વપરાતી અનૌપચારિક
તર્કપ્રક્રિયાને નજીકથી અનુસરવા માટે રચાયાં છે (એટલે તે કંઈક અંશે
``સ્વાભાવિક'' છે). સ્વાભાવિક નિગમનની સાબિતીઓ ધારણાઓથી શરૂ થાય છે.
પછી અનુમાન-નિયમો લાગુ પાડવામાં આવે છે. \Intro{\lnot}, \Intro{\lif},
\iftag{FOL}{\Elim{\lor} અને \Elim{\lexists}}{ તથા \Elim{\lor}}
અનુમાન-નિયમો ધારણાઓને ``!!{discharged}'' કરે છે; સ્પષ્ટતા માટે
!!{discharged} ધારણાનું ચિહ્ન અનુમાનની બાજુમાં મૂકવામાં આવે છે.
\end{explain}

\begin{defn}[ધારણા]
કોઈ પણ શાખામાં સૌથી ઉપરના સ્થાને આવતું કોઈ પણ !!{sentence}
\emph{ધારણા} છે.
\end{defn}

સ્વાભાવિક નિગમનમાં !!^{derivation}s એ !!{sentence}sનાં અમુક વૃક્ષો છે.
તેમાં સૌથી ઉપરનાં !!{sentence}s ધારણાઓ હોય છે; અને જો !!a{sentence}
બીજા એક, બે કે ત્રણ !!{sentence}sની નીચે આવે, તો તે અનુમાનના કોઈ નિયમથી
યોગ્ય રીતે ફલિત થવું જોઈએ. અનુમાનની ઉપરનાં !!{sentence}sને
\emph{આધારવાક્યો} અને નીચેના !!{sentence}ને અનુમાનનું \emph{ફલિતવાક્ય}
કહે છે. નિયમો જોડીઓમાં આવે છે: દરેક !!{operator} માટે એક પરિચય-નિયમ
અને એક નિવારણ-નિયમ. તેઓ ફલિતવાક્યમાં !!a{operator} દાખલ કરે છે અથવા
નિયમના કોઈ આધારવાક્યમાંથી !!a{operator} દૂર કરે છે. કેટલાક નિયમો અમુક
પ્રકારની ધારણાને \emph{!!{discharged}} કરવાની છૂટ આપે છે. કઈ ધારણા કયા
અનુમાનથી !!{discharged} થાય છે તે બતાવવા માટે ધારણા અને અનુમાન બંનેને
ચિહ્ન આપીએ છીએ. ધારણાને ``$\Discharge{!A}{n}$'' એમ લખીને આ દર્શાવાય છે.
\sourcecorrection{OLND-001}{સ્થિર અંગ્રેજી મૂળમાં આ વૃક્ષવર્ણનના એક સ્થાને
``બીજા સિક્વન્ટો'' લખાયું છે; સ્વાભાવિક નિગમનનું વૃક્ષ વાક્યોનું હોવાથી
અહીં આશય મુજબ ``બીજાં વાક્યો'' વાંચ્યું છે.}

% Only include this sentence is on of \land, lor, lif or lnot is defined.

\iftag{notprvNot,notprvAnd,notprvOr,notprvIf}{}%
	{$\land$, $\lor$, $\lif$, $\lnot$ અને $\lfalse$માંથી કેટલાક
!!{operator}s વ્યાખ્યાયિત હોય તોપણ, તે બધા માટે નિયમો ધ્યાનમાં લેવાનો
રિવાજ છે.}



\end{document}
